{\noindent\large\bfseries Abstract}\\[0.4em]
{\noindent\small
Three companion papers established that the area-weighted mean surface gravity
$\gmean$ of a small body can be estimated from a polyhedral shape model,
without an independently determined density (Papers~I--III). A single scalar,
however, cannot say \emph{where} on a body the gravity is high or low, how
steeply the surface is inclined to the local effective gravity, or where loose
regolith is driven to accumulate. Here we extend the shape-based framework from
the scalar mean to the spatially resolved surface fields: the effective gravity
vector $\geff$ (self-gravitation plus rotation), the gravitational slope
$\slope$, and the surface geopotential $\Upot$, evaluated facet by facet for
seven asteroids from a polyhedral shape model together with an assumed uniform
density and a specified spin state.

The framework is validated against an independent per-facet reference. For
(162173)~Ryugu we compare with the gravitational products released by the
Hayabusa2 team on the identical $49{,}152$-facet shape model at the same
density and rotation period. Agreement is essentially exact: mean slope
$15.06^{\circ}$ against $15.06^{\circ}$, RMS difference $0.48^{\circ}$,
$99.96\%$ of facets within $5^{\circ}$ and $99.95\%$ within $2^{\circ}$,
Pearson $r=0.9989$ in slope and $0.9898$ in geopotential. For (101955)~Bennu the global
mean slope is $15.5^{\circ}$ and the latitudinal profile rises from
$10.5^{\circ}$ near the equator to $19.9^{\circ}$ at mid-latitude before
falling toward the pole, reproducing the pattern reported from OSIRIS-REx data.

The principal methodological result concerns stability under mesh refinement.
Because slope is a \emph{local} quantity, defined through individual facet
normals, the area-weighted mean slope varies by $1.8\%$ to $25.6\%$ across
meshes of the same body decimated from roughly $42{,}000$ to $2{,}000$ facets,
whereas the normalized geopotential dispersion $\potvarn$ varies by only $0.9\%$
to $5.9\%$, and is the more stable of the two on every body tested. The
sharper distinction is stabilization: between the two finest meshes the
geopotential dispersion is already near a plateau (changes of $0.00\%$ to
$0.60\%$), whereas the mean slope shows no comparable plateau within the tested
resolution range, still rising by up to $8.2\%$. A slope value is therefore not
reproducible unless the mesh resolution is quoted with it, whereas the
geopotential dispersion exhibits practical convergence over the tested range.

Comparing two independent shape models of Itokawa---the spacecraft-derived
Gaskell model and a ground-based radar model---separates what survives a coarser
shape from what does not. Bulk descriptors are compared at a common facet count
by interpolation, and distribution statistics at the nearest available direct
meshes. Bulk descriptors agree
closely (area-weighted mean slope within $0.8\%$, geopotential dispersion within
$7.9\%$), but the distribution tails do not: the radar model reaches only
$41^{\circ}$ with $1.7\%$ of facets above $30^{\circ}$, against $149^{\circ}$
and $4.8\%$ for the Gaskell model. In this controlled case a ground-based shape
therefore preserves the global gravity descriptors while failing to reproduce
the distribution tails; one case does not establish the result generally, but it
indicates where coarse models may and may not be relied upon. We verify the computed fields against two exact solutions: the
closed-form potential of a homogeneous ellipsoid (Pearson $r=1.000000$, mean
relative error $0.18\%$) and the analytic slope field of a rotating sphere.
The latter test is essential---a non-rotating sphere check cannot detect an
error in the centrifugal term---and we recommend it as a standard verification
for surface-gravity codes. Both tests are released as runnable scripts together
with the analysis code.
}

\medskip
\noindent\textbf{Keywords:} Asteroids; Surface gravity; Gravitational slope;
Geopotential; Shape models; Polyhedral gravity

\section{Introduction}
\label{sec:intro}

The surface gravity of a small body governs nearly every process that shapes
its regolith and constrains how a spacecraft can operate on or near it. Where
material comes to rest, how loose regolith migrates, whether a boulder is
stable on a slope, and where a lander or sampling head can safely touch down
are all set by the local effective gravity---the vector sum of self-gravitation
and the centrifugal acceleration of rotation
\citep{scheeres2016,scheeres2019}. For the bodies visited by spacecraft these
quantities are now mapped in detail from mission data. For the far larger
population resolved only as a shape model---from ground-based radar, lightcurve
inversion, or a single flyby---no such maps exist, even though the same physics
applies.

The companion papers of this series developed a purely shape-based route to the
\emph{mean} surface gravity. Paper~I \citep{paper1} showed that the
area-equivalent radius $\RA$ outperforms the volume-equivalent radius $\RV$ for
elongated bodies and proved the bound $GM/\RA^{2}\leq\gmean$; Paper~II
\citep{paper2} derived the optimal blended radius from a non-analytic
$\Pthree/\Ptwo$ deformation law; and Paper~III \citep{paper3} tested that
estimator against a homogeneous-polyhedron reference field for $22$ non-convex
models, obtaining a mean absolute error of $1.72\%$, and delimited its domain of
validity. That work establishes that one scalar---the area-weighted mean
gravity---can be recovered from the shape model to within a few per cent. A
scalar
mean, however, is silent about the spatial structure of the field. Those
questions require moving from a single number to a field defined over the
entire surface.

This paper makes that step. Using the same polyhedral shape models and an
assumed uniform density, we compute three spatially resolved quantities at
every facet: the effective gravity $\geff$, the gravitational slope $\slope$
(the angle between $\geff$ and the outward facet normal $\nhat$), and the
geopotential $\Upot$, whose minima mark the sinks toward which loose material
migrates. The self-gravitation is evaluated with the analytic polyhedral method
of \citet{werner1996}, as implemented in a modern open-source library
\citep{tsoulis2012,gravitylib}, so that the framework remains light enough to
run on a laptop yet remains accurate on the concave, bilobed, and highly
irregular shapes that dominate the small-body population.

Making these fields easy to compute from a shape model, an assumed density and
a specified spin state is useful, but it also
exposes a methodological hazard that, to our knowledge, has not been stated
plainly in the shape-based context: the slope field is sensitive to the
resolution of the shape model. Because slope is a local quantity, defined
through the orientation of each individual facet, refining a mesh changes the
distribution of facet normals and hence the mean slope, even when the underlying
body is identical. Section~\ref{sec:slopesens} shows that over a
$2{,}000$-to-$42{,}000$-facet range the area-weighted mean slope varies by up to
a quarter of its own value, and---more tellingly---that it has still not
converged at the finest meshes available. A slope value quoted without its mesh
resolution is therefore not reproducible, a point with direct consequences for
how such numbers are reported and compared between studies.

The central positive result of this paper is that a closely related
quantity---the area-weighted dispersion of the surface geopotential---behaves far
better. Because the geopotential is an integral of the mass distribution rather
than a local surface property, the facet-scale roughness that scatters the slope
largely cancels in the potential. Across the bodies tested, $\potvarn$ is the
more stable of the two diagnostics in every case, and it approaches a practical
plateau over the tested refinement range whereas the slope shows no comparable
stabilization (Section~\ref{sec:varrobust}). This is
consistent with the behaviour noted by \citet{kanamaru2019} for Itokawa and
makes $\potvarn$ the natural descriptor to quote when a resolution-independent
statement about a body's surface-gravity field is required.

A separate question is what survives when the shape model itself is coarse, as
it must be for objects never visited by a spacecraft. Itokawa is the rare case
with two independent shape models of comparable extent---the spacecraft-derived
Gaskell model and a ground-based radar model---and comparing them at matched
resolution (Section~\ref{sec:crossversion}) separates the descriptors that
tolerate a ground-based shape from those that do not.

We validate the framework against bodies for which detailed gravitational
products already exist. For (162173)~Ryugu we compare directly against the
per-facet gravitational products released with \citet{kanamaru2019}, computed on
the identical shape model at the same density and rotation period; this is the
strongest available test, since shape, method, density, and spin are all held
fixed. For (101955)~Bennu the recovered mean slope and latitudinal profile are
compared with the OSIRIS-REx results
\citep{scheeres2019,barnouin2020,jawin2020}. We then use (433)~Eros,
(25143)~Itokawa, (951)~Gaspra, (243)~Ida, and (216)~Kleopatra to characterise
the resolution behaviour of the two diagnostics across the range of small-body
morphologies \citep{zuber2000,thomas2002,fujiwara2006,abe2006}.

\medskip
Three methodological commitments govern the analysis.

First, \emph{every slope figure is reported with its mesh resolution}. Because
the sensitivity documented here is a property of the diagnostic rather than of
the body, a slope value quoted without a facet count is not reproducible. All
slope numbers in this paper are accompanied by the facet count at which they
were obtained. Cross-model comparisons are made at matched resolution wherever
the meshes permit it, and where they do not---as for the tails of the Itokawa
distributions---the facet counts actually used are stated.

Second, \emph{the fields are verified against exact solutions, including one
that exercises the rotating frame}. The polyhedral potential is compared with
the closed-form potential of a homogeneous ellipsoid, and the full slope field
with the analytic slope of a uniformly rotating sphere. The second test is not
redundant: a non-rotating sphere reproduces zero slope for either sign of the
centrifugal term, so it cannot detect an error there, whereas the rotating case
discriminates sharply. We report both, and recommend the rotating test as a
standard check for codes of this kind.

Third, \emph{the uniform-density assumption is stated rather than hidden}. All
fields in this paper are computed for a homogeneous interior. This assumption
holds exactly for none of the bodies considered, and the resulting
limitations---in particular for Itokawa, whose head and body are known to
differ in density---are examined in Section~\ref{sec:discussion} rather than
deferred.

The paper is organised as follows. Section~\ref{sec:methods} defines the
effective gravity, slope, and geopotential, describes the polyhedral solver and
the shape models used, and specifies the decimation and smoothing operators.
Section~\ref{sec:validation} presents the validation against Ryugu and Bennu.
Section~\ref{sec:slopesens} documents the resolution behaviour of the slope,
and Section~\ref{sec:varrobust} that of the geopotential dispersion.
Section~\ref{sec:crossversion} compares independent shape versions of the same
body. Section~\ref{sec:discussion} discusses the consequences for cross-study
comparison and for interior-structure inference, together with the limitations
of the uniform-density assumption, and Section~\ref{sec:conclusions} concludes.
